% Part: history
% Chapter: biographies
% Section: ernst-zermelo

\documentclass[../../../include/open-logic-section]{subfiles}

\begin{document}

\olfileid{his}{bio}{zer}

\olsection{ఎర్న్‌స్ట్ సెర్మెలో}

\olphoto{zermelo-ernst}{ఎర్న్‌స్ట్ సెర్మెలో}

ఎర్న్‌స్ట్ సెర్మెలో 1871 జూలై 27న జర్మనీలోని బెర్లిన్‌లో జన్మించారు. ఆయనకు ఐదుగురు అక్కాచెల్లెళ్లు ఉన్నారు; కానీ కుటుంబంలో అనారోగ్య సమస్యల వల్ల వారిలో ముగ్గురే పెద్దవయసు వరకు జీవించారు. ఆయన చిన్నతనంలోనే తల్లిదండ్రులు కూడా మరణించారు; పదిహేడేళ్ల వయసుకే ఆయన, ఆయన తోబుట్టువులు అనాథలయ్యారు. సెర్మెలోకు కళలపై, ముఖ్యంగా కవిత్వంపై, గాఢ ఆసక్తి ఉండేది. ఆయన చురుకైన, చమత్కారమైన, విమర్శనాత్మక స్వభావానికి పేరుగాంచారు. ఆయన అత్యంత ప్రసిద్ధ గణిత విజయాల్లో ఎంపిక స్వయంసిద్ధాన్ని ప్రవేశపెట్టడం (1904), సమితి సిద్ధాంతాన్ని స్వయంసిద్ధీకరించడం (1908) ఉన్నాయి.

విశ్వవిద్యాలయంలో సెర్మెలో ఆసక్తులు విస్తృతంగా ఉండేవి. భౌతికశాస్త్రం, గణితం, తత్వశాస్త్రంలో పాఠ్యాంశాలు చదివారు. హెర్మాన్ శ్వార్జ్ పర్యవేక్షణలో బెర్లిన్ విశ్వవిద్యాలయంలో \emph{Investigations in
  the Calculus of Variations} అనే పరిశోధన గ్రంథాన్ని 1894లో పూర్తి చేశారు. 1897లో మరింత చదువు కోసం గ్యోటింగెన్ విశ్వవిద్యాలయానికి వెళ్లాలని నిర్ణయించుకున్నారు; అక్కడ డేవిడ్ హిల్బర్ట్ పునాదుల సంబంధిత కృషి ఆయనపై బలమైన ప్రభావం చూపింది. 1899లో ఆచార్య పదవికి అర్హత సాధించినా, పదకొండు సంవత్సరాల తరువాతే ఆ పదవి దక్కింది---దీనికి ఆయన విచిత్ర ప్రవర్తన, ``నాడీ ఆతురత'' కారణాలై ఉండవచ్చు.

చివరకు 1910లో జ్యూరిక్ విశ్వవిద్యాలయంలో సెర్మెలోకు జీతంతో కూడిన ఆచార్య పదవి లభించింది; కానీ క్షయవ్యాధి కారణంగా 1916లో పదవీ విరమణ చేయవలసి వచ్చింది. కోలుకున్న తరువాత 1921లో ఫ్రైబర్గ్ విశ్వవిద్యాలయంలో గౌరవ ఆచార్య పదవి ఇచ్చారు. ఈ కాలంలో గణిత పునాదులపై పనిచేశారు. థోరాల్ఫ్ స్కోలెమ్, కుర్ట్ గ్యోడెల్ కృషిపై ఆయన అసంతృప్తి చెందారు; తన వ్యాసాల్లో వారి పద్ధతులను బహిరంగంగా విమర్శించారు. ప్రజాదరణ లేకపోవడం, జర్మనీలో హిట్లర్ అధికారంలోకి రావడాన్ని వ్యతిరేకించడం వల్ల 1935లో ఫ్రైబర్గ్ పదవి నుంచి తొలగించారు.
 
సెర్మెలో జీవితపు చివరి సంవత్సరాలు ఒంటరితనంతో గడిచాయి. 1935లో తొలగింపు తరువాత గణితాన్ని వదిలేశారు. గ్రామీణ ప్రాంతానికి మారి సాధారణ జీవితం గడిపారు. 1944లో వివాహం చేసుకున్నారు; చూపు తగ్గుతుండటంతో భార్యపై పూర్తిగా ఆధారపడేవారు. 1951 నాటికి సెర్మెలో చూపును పూర్తిగా కోల్పోయారు. 1953 మే 21న జర్మనీలోని గ్యున్టర్‌స్టాల్‌లో మరణించారు.

\begin{reading}
సెర్మెలో సమగ్ర జీవిత చరిత్రకు \citet{Ebbinghaus2015} చూడండి. 1904, 1908 సంవత్సరాల ఆయన ముఖ్య వ్యాసాలు మూల జర్మన్‌లో చదవడానికి లభిస్తాయి \citep{Zermelo1904,Zermelo1908}. భౌతికశాస్త్రంపై రచనలతో సహా సెర్మెలో సంకలిత రచనలు ఆంగ్ల అనువాదంలో \citep{Ebbinghaus2010,Ebbinghaus2013}లో లభిస్తాయి.
\end{reading}

\end{document}
